%% Entropy (MDPI) — CGWT: Computational Constraints on Coalition Broadcast and
%% Causal Integrated Information in a Recurrent Neural Processor System
%% https://www.mdpi.com/journal/entropy

\documentclass[entropy,article,accept,moreauthors,pdftex]{Definitions/mdpi}

\firstpage{1}
\makeatletter
\setcounter{page}{\@firstpage}
\makeatother
\pubvolume{27}
\issuenum{1}
\articlenumber{0}
\pubyear{2026}
\copyrightyear{2026}

\Title{What Coalition Broadcast Reveals About Causal Integration:\\
A Diagnostic Toolkit and Empirical Constraints from a GWT--IIT Neural Processor System}

\Author{Guyu Adam$^{1}$}

\AuthorNames{Guyu Adam}

\address{%
$^{1}$ \quad Independent Research; dhuojian@gmail.com}

\corres{Correspondence: dhuojian@gmail.com}

%══════════════════════════════════════════════════════════════════════════════
% Abstract
%══════════════════════════════════════════════════════════════════════════════

\abstract{%
Global Workspace Theory (GWT) and Integrated Information Theory (IIT) offer
complementary but mechanistically distinct accounts of conscious access,
yet whether GWT-style broadcast architectures can produce the high-$\Phi$
states that IIT identifies as the hallmark of consciousness remains an open
computational question. Rather than testing a specific theoretical prediction,
we construct a \textit{diagnostic testbed}—a five-processor recurrent neural
network system (512 neurons each, 2560 total) with real causal structure—and
systematically measure what a tractable $\Phi$ approximation actually captures
across nine experimental modes (900 observation cycles, GPU-accelerated).
We find that among multi-processor modes, $\Phi_\text{approx}$ falls in a
narrow band of 1.022--1.049 regardless of broadcast mechanism. $\Phi$ scales
approximately linearly with total recurrent neuron count ($\sim$2.2$\times$
for five processors vs.\ one), and its inter-processor integration component
($\Phi_\text{between}$) is a near-constant 0.05 across all conditions.
Processor differentiation provides a small but statistically reliable benefit
(clustered Cohen's $d = 0.13$, $p < 0.001$). Coalition consensus broadcast,
after correcting a pairwise agreement mechanism that prevented multi-member
coalition formation in the initial implementation, successfully forms
coalitions (effective ensemble size $\approx 2.5$) but confers no $\Phi$
advantage over single-winner selection. Two implementation issues—an effective
information accumulation artifact and the coalition formation bug—were
identified \textit{through} the diagnostic tools developed in this work
($\Phi$ decomposition, attention weight tracking, per-cycle monitoring).
The primary contribution is methodological: a transferable diagnostic toolkit
for computational studies of broadcast-mediated information integration,
together with empirical constraints on what $\Phi_\text{approx}$ measures in
recurrent architectures.
}

\keyword{%
global workspace theory; integrated information theory; phi; coalition broadcast;
recurrent neural network; causal integration; multi-agent coordination}

%══════════════════════════════════════════════════════════════════════════════
% 1. Introduction
%══════════════════════════════════════════════════════════════════════════════

\begin{article}

\section{Introduction}

Two theoretical frameworks dominate the contemporary scientific study of
consciousness. Global Workspace Theory (GWT)~\cite{baars1988cognitive,
dehaene2011experimental} proposes that conscious access occurs when information
from specialized processors wins a competition for broadcast into a
limited-capacity global workspace---a ``winner-take-all'' architecture.
Integrated Information Theory (IIT)~\cite{tononi2004information,tononi2016integrated}
proposes that consciousness is identical to the amount of irreducible
integrated information ($\Phi$) generated by a system's causal
structure---how much more information the whole system specifies about its past
and future states than the sum of its parts.

These two frameworks make complementary but mechanistically distinct
predictions. GWT emphasizes the \textit{selection} of information for global
access; IIT emphasizes the \textit{integration} of information into an
irreducible whole. Whether GWT's competitive broadcast architecture can
produce the high-$\Phi$ states that IIT identifies as the hallmark of
consciousness remains an open question---and one that has recently been
subjected to adversarial empirical testing~\cite{ferrante2025adversarial}.

The tension is sharpened by four recent developments. First,
Ferrante et al.\ (2025)~\cite{ferrante2025adversarial} demonstrated, through
a preregistered adversarial collaboration, that both GNWT and IIT receive
only partial empirical support---establishing rigorous theory-testing through
controlled computational and experimental manipulation as the methodological
standard for the field. Second,
Luppi et al.\ (2024)~\cite{luppi2024synergistic} demonstrated that the human
brain's global workspace is characterized by a specific balance of synergistic
and redundant information, raising the question of whether artificial GWT
architectures can reproduce this balance. Third,
Kearney (2026)~\cite{kearney2026information} established a formal bridge
between IIT and the Free Energy Principle via Maximum Caliber, providing a
theoretical framework for linking $\Phi$ to predictive processing. Fourth,
Barrett et al.\ (2026, preprint, not peer-reviewed)~\cite{barrett2026iit}
clarified that $\Phi$ cannot be precisely defined for real neural systems,
necessitating principled approximation strategies for any computational
implementation.

In parallel, the multi-agent AI literature has independently converged on
questions about broadcast-based coordination among LLM-based
agents~\cite{multiagent_coordination}. However, these works typically measure
task performance rather than causal information integration, and the dominant
debate paradigm conflates broadcast architecture with agent implementation.

\subsection{A Diagnostic Testbed for Broadcast-Mediated Integration}

Rather than testing a specific theoretical prediction, we construct a
\textit{diagnostic testbed}---a minimal recurrent neural processor system
in which broadcast architecture can be systematically varied while measuring
$\Phi$ from the same underlying causal structure. The testbed implements
nine broadcast modes (competitive, random, no-broadcast, collaborative,
hybrid, adaptive, single-processor, single-self-broadcast, and homogeneous
competitive), each differing only in how information is selected for global
workspace access. By holding the processor architecture constant and varying
only the selection mechanism, we isolate the causal contribution of broadcast
to $\Phi$.

Collaborative Global Workspace Theory (CGWT)---which replaces winner-take-all
selection with coalition consensus broadcast---serves as one of the broadcast
architectures under test, not as the hypothesis to be confirmed. Prior
GWT--IIT integrations~\cite{iwmt2020, safron2020integrated} and adjacent
implementations~\cite{godelos} have proposed coalition-like mechanisms, but
none have systematically compared coalition against competitive, random, and
null-broadcast baselines on a shared causal substrate.

Our approach differs from prior work in three ways: (1) we use recurrent
neural network (RNN) processors with real causal structure, enabling
$\Phi$ to be computed from neural state transition probability matrices
rather than token-distribution proxies (LLM-based) or correlation patterns
(fMRI-based); (2) we deploy a diagnostic toolkit---$\Phi$ decomposition
into within- and between-processor components, attention weight tracking,
per-cycle $\Phi$ trajectory monitoring, and coalition merge fidelity
metrics---that enables not just hypothesis testing but \textit{diagnosis of
why} a given broadcast mechanism does or does not affect $\Phi$; and (3) we
compare nine modes head-to-head on identical stimuli, providing a
decomposition of the factors contributing to $\Phi$ in recurrent
architectures.

Seven pre-registered hypotheses (H1--H7, detailed in Methods) guided the
experimental design. We report all results regardless of outcome, accepting
null results where the data warrant them. The contribution is not the
confirmation or disconfirmation of any specific hypothesis, but the
diagnostic framework and the empirical constraints it reveals about what
$\Phi_\text{approx}$ measures in recurrent neural architectures.

\subsection{Contributions and Positioning}

The present work makes four contributions:

\begin{enumerate}
\item \textbf{A diagnostic toolkit for computational $\Phi$ studies.}
The $\Phi_\text{within}$ / $\Phi_\text{between}$ decomposition, attention
weight distribution tracking, coalition size monitoring, per-cycle $\Phi$
trajectory analysis, and merge fidelity metrics constitute a transferable
toolkit. These tools identified two non-trivial implementation issues during
development (an effective information accumulation artifact from expanding
TPM windows and a coalition formation bug from mismatched agreement metrics),
demonstrating their diagnostic value independently of any specific empirical
finding.

\item \textbf{A causal $\Phi$ pipeline from RNN state transitions.}
We establish a complete computational pipeline---RNN activation $\to$
state clustering $\to$ causal TPM $\to$ effective information /
mutual information $\to$ composite $\Phi$ approximation---that respects
the causal structure requirement of IIT while remaining tractable at scale.

\item \textbf{Empirical constraints on $\Phi_\text{approx}$ in recurrent
architectures.} Across 900 observation cycles, we find that (a)
$\Phi_\text{approx}$ is dominated by total recurrent neuron count, not by
broadcast mechanism; (b) processor differentiation provides a small but
reliable $\Phi$ benefit (clustered $d = 0.13$); and (c) coalition consensus
broadcast, even when correctly implemented, confers no $\Phi$ advantage over
single-winner selection under the current $\Phi$ approximation.

\item \textbf{Operationalizing the MaxCal IIT--FEP bridge.}
We translate the Maximum Caliber bridge of
Kearney~(2026)~\cite{kearney2026information} into a concrete coalition
selection objective (Eq.~\ref{eq:maxcal_coalition}), demonstrating how a
mathematical framework can become a design principle for multi-agent
architectures.
\end{enumerate}

This work does \textit{not} claim that the RNN system is conscious, that
$\Phi_\text{approx}$ equals true $\Phi$, or that any specific broadcast
architecture is superior. Its contribution is the diagnostic methodology
and the empirical constraints it reveals---a computational contribution to
the methodology of theory-testing through precise, falsifiable predictions
advocated by Ferrante et al.\ (2025) and Phua~(2025)~\cite{phua2025}.

%══════════════════════════════════════════════════════════════════════════════
% 2. Methods
%══════════════════════════════════════════════════════════════════════════════

\section{Methods}

\subsection{System Architecture}

Our system consists of five recurrent neural network (RNN) processors
connected to a global workspace (Figure~\ref{fig:architecture}). Each processor
is a small RNN with $N = 512$ tanh neurons and a stimulus input of dimension
$d_\text{in} = 32$. The processor dynamics are:

\begin{equation}
h_t = \tanh(W_\text{in} \cdot s + W_\text{rec} \cdot h_{t-1} + b + \eta_t)
\label{eq:rnn}
\end{equation}

where $s \in \mathbb{R}^{32}$ is the encoded stimulus vector,
$W_\text{rec} \in \mathbb{R}^{512 \times 512}$ is the recurrent weight matrix,
$W_\text{in} \in \mathbb{R}^{512 \times 32}$ projects the stimulus into the
recurrent space, $b$ is a bias term, and $\eta_t \sim \mathcal{N}(0, 0.01)$ is
Gaussian noise added at each of $T = 25$ unroll steps. The processor's
``proposal'' is its final hidden state $h_T$.

The five processors are structurally differentiated through distinct
$W_\text{rec}$ initialization patterns (Table~\ref{tab:processors}),
producing qualitatively different dynamical regimes. All initializations are
parameterization-independent (scale to arbitrary $N$) and enforce the
echo-state property (spectral radius $< 0.9$) for stable dynamics.

\begin{table}[h]
\caption{Five specialized processor types and their recurrent connectivity
patterns. All enforce $\rho(W_\text{rec}) < 0.9$.}
\label{tab:processors}
\centering
\begin{tabular}{lll}
\toprule
Processor & Connectivity & Dynamical Signature \\
\midrule
Perceptor & Sparse near-diagonal, local lateral & Fast decorrelation, feature extraction \\
Reasoner  & Chain-structured, forward-biased & Sequential processing stages \\
Evaluator & Two-pool bistable attractor & Decision-making, winner-take-all dynamics \\
Integrator & Watts--Strogatz small-world & Holistic integration, long-range coupling \\
Predictor & Upper-triangular dominant & Anticipatory, temporal prediction \\
\bottomrule
\end{tabular}
\end{table}

The global workspace stores the most recent broadcast content. Processors
receive the workspace state as a context vector injected into their hidden
state before processing the next stimulus: $h_0' = 0.4 \cdot h_0 + 0.6 \cdot
\text{ws}$, where $\text{ws}$ is the broadcast activation vector from the
previous cycle.

\subsection{Experimental Modes}

We define nine experimental modes (Table~\ref{tab:modes}). Modes 1--6 are
core conditions testing the CGWT hypothesis; modes 7--9 are methodological
controls (self-broadcast, homogeneous, adaptive gate).

\begin{table}[h]
\caption{Nine experimental modes. Modes 7--9 are methodological controls.}
\label{tab:modes}
\centering
\begin{tabular}{llp{5cm}}
\toprule
Mode & Selection Mechanism & Scientific Purpose \\
\midrule
1. Competitive    & Winner-take-all via attention scoring   & GWT baseline \\
2. Random         & Uniform random winner selection         & Null selection baseline \\
3. No-broadcast   & No broadcast (processors isolated)      & Null broadcast baseline \\
4. Collaborative  & Coalition consensus (2 members)         & CGWT core prediction \\
5. Hybrid         & Top-2 competitive $\to$ coalition merge & CGWT with competitive preselection \\
6. Single-processor & One processor, no competition        & Multi-processor necessity test \\
7. Single self-broadcast & One processor broadcasts to itself & Isolate broadcast vs.\ multi-processor \\
8. Homogeneous competitive & Five identical Integrators, winner-take-all & Processor differentiation necessity \\
9. Adaptive & Coalition with differentiation gate (std $<$ 0.02 $\to$ competitive fallback) & Test adaptive gating$^a$ \\
\bottomrule
\end{tabular}

\smallskip
\small
$^a$ The threshold $\tau = 0.02$ was chosen to be approximately one order
of magnitude above the observed attention weight noise floor
($\sigma(\alpha) \approx 0.001$--$0.003$ in collaborative/hybrid modes),
such that coalition is deployed only when processor disagreement exceeds
what can be attributed to numerical variation. \\
\end{table}

\subsection{$\Phi$ Approximation}

Computing IIT's exact $\Phi$ as defined in IIT 3.0~\cite{oizumi2014}
for a system of $5 \times 512 = 2560$ neurons is
cosmologically intractable (state space size $2^{2560}$). We compute a
tractable approximation with four documented simplifications:

\begin{enumerate}
\item \textbf{State space reduction.} Continuous activation vectors are
clustered into $k = 50$ macro-states via $k$-means (5 random restarts,
fixed seed for reproducibility). The state transition probability matrix
(TPM) is built from the sequence of cluster assignments.

\item \textbf{No MIP search.} IIT's $\Phi$ requires finding the Minimum
Information Partition---the bipartition that minimizes information loss.
This is NP-hard. We use a fixed reduction: $\Phi \approx I(\text{workspace};
\text{all}) - \frac{1}{n}\sum_i I(\text{workspace}; \text{processor}_i)$.
Random partition sampling (100 bipartitions per cycle) confirmed that the
bias from this fixed partition is consistent across conditions ($\sigma_\Phi <
0.3 \cdot \mu_\Phi$).

\item \textbf{Histogram-based MI.} Mutual information is estimated via 2D
histogram with quantile (equi-frequency) binning, making no distributional
assumptions. The spatial-ergodicity assumption (neuron indices as paired
observations) follows standard practice in representational similarity
analysis~\cite{kriegeskorte2008representational}. Adaptive bin count ($\min(20,
n/4)$) prevents over-binning of small samples.

\item \textbf{Weighted composite.} $\Phi_\text{approx} = \omega_\text{mi}
\cdot \text{MI}_\text{integration} + \omega_\text{ei} \cdot \text{EI} +
\omega_\text{diff} \cdot \text{differentiation}$, with default weights
$(0.40, 0.35, 0.25)$. Sensitivity analysis using
Dirichlet$([4, 3.5, 2.5])$ sampling (500 draws, concentrated around default
weights) produces coefficients of variation (CV) of 0.35--0.51 across modes.
While these exceed our conservative $< 0.3$ threshold, all modes occupy a
narrow CV range, indicating that the estimation bias is consistent across
conditions and does not threaten comparative rank-ordering. We recommend
permutation-based CV baselines for future work to replace the fixed threshold.
\end{enumerate}

\textbf{Critical honesty.} This is not IIT's proper $\Phi$. The absolute
values are meaningless; only the comparative rank ordering across experimental
conditions carries scientific weight. Validation of $\Phi_\text{approx}$
against exact $\Phi$ on small tractable systems (e.g., $N = 6$--$12$ binary
units) remains as future work.

\subsection{$\Phi$ Decomposition}

To isolate the contributions of processor-internal recurrence versus
processor-to-processor communication, we decompose $\Phi$ into two components:

\begin{equation}
\Phi_\text{total} = \Phi_\text{within} + \Phi_\text{between}
\label{eq:decompose}
\end{equation}

$\Phi_\text{within}$ is the average effective information (EI) computed from
each processor's individual activation TPM---capturing the causal structure
intrinsic to the recurrent weights $W_\text{rec}$. $\Phi_\text{between}$ is
the MI integration component: $I(\text{workspace}; \text{all processors}) -
\frac{1}{n}\sum_i I(\text{workspace}; \text{processor}_i)$, weighted by the
$\omega_\text{mi}$ and $\omega_\text{diff}$ coefficients. The decomposition
shares the same weighted-composite structure as $\Phi_\text{approx}$ but
separates the processor-internal EI contribution. Note that
$\Phi_\text{within}$ and $\Phi_\text{after}$ are computed from different
algorithmic paths and have different scales; only cross-mode comparisons of
each component (not their ratios) are meaningful.

\subsection{Coalition Merge Strategy}

In collaborative and hybrid modes, the coalition's outputs are merged into a
single broadcast representation using attention-weighted averaging:

\begin{equation}
\text{merged} = \sum_{i \in \text{coalition}} \alpha_i \cdot \vec{v}_i,
\quad \alpha_i = \text{softmax}(\cos(\vec{v}_i, \vec{ws}))
\label{eq:merge}
\end{equation}

where $\vec{v}_i$ is processor $i$'s activation vector and $\vec{ws}$ is the
current workspace state. This preserves differentiation better than a simple
mean (which cancels complementary information), but its effectiveness depends
on processors producing sufficiently differentiated representations.

\subsection{Experimental Procedure}

Each experimental cycle consists of:

\begin{enumerate}
\item \textbf{Pre-broadcast $\Phi$ measurement.} $\Phi$ is computed from the
current workspace state, processor proposals, and workspace history.
\item \textbf{Stimulus processing.} All processors process the same stimulus
vector through their respective RNNs.
\item \textbf{Selection (mode-dependent).} Competitive: attention-based
winner-take-all. Collaborative: consensus-based coalition selection (2
members). Hybrid: competitive preselection to top-2, then coalition merge.
Random: uniform random selection. No-broadcast: skip.
\item \textbf{Broadcast.} The winner (or merged coalition representation)
is broadcast into the global workspace and injected into all processors
for the next cycle.
\item \textbf{Post-broadcast $\Phi$ measurement.} $\Phi$ is recomputed after
the broadcast.
\end{enumerate}

\subsection{Stimuli and Data Collection}

Twenty stimuli spanning consciousness theory and multi-agent coordination were
encoded into 32-dimensional vectors via deterministic hashing with random
projection (same text always produces the same vector). Each stimulus was
processed for 5 cycles per mode, yielding $20 \times 5 = 100$ observations per
mode and $100 \times 9 = 900$ total observation cycles. Stimuli included
questions such as ``What is the relationship between perception and
consciousness?'' and ``How should an agent weight novel information versus
prior consensus?''

Experiments ran on an NVIDIA RTX 5060 Ti (16 GB VRAM) with CuPy GPU
acceleration. Each cycle completes in approximately 0.1 seconds. All code,
data, and experiment configurations (including dependency versions) are
archived with the manuscript.

\subsection{Pre-Registered Hypotheses}

The experimental design was guided by seven pre-registered hypotheses,
reported here for transparency. All results are reported regardless of
outcome. The dependent variable (phi\_after for absolute level; phi\_delta
for broadcast gain) is declared per hypothesis based on the scientific
question each addresses.

\textbf{Phi-after hypotheses (absolute level):} H5b---broadcast modes produce
substantially higher $\Phi$ than no-broadcast. H7b---structurally
differentiated processors produce higher $\Phi$ than homogeneous processors.

\textbf{Phi-delta hypotheses (broadcast gain):} H1--H4---coalition modes
(collaborative, hybrid) produce larger $\Phi$ increases than competitive or
random selection. H5a---competitive broadcast produces larger $\Phi$ increases
than no-broadcast. H6---single-processor self-broadcast produces equivalent
$\Phi$ to single-processor without broadcast. H7a---differentiated processors
produce larger $\Phi$ increases than homogeneous processors.

\subsection{Statistical Analysis}

We report two complementary statistical approaches:

\textbf{Non-parametric pairwise tests.} Mann--Whitney $U$ tests (one-tailed,
$\alpha = 0.05$) with Cohen's $d$ effect sizes and required sample sizes for
80\% power. The dependent variable (phi\_after or phi\_delta) is
declared per hypothesis based on the scientific question: phi\_after for
hypotheses about absolute levels (H5b, H7b), phi\_delta for hypotheses about
broadcast gain (H1--H4, H5a, H6, H7a).

\textbf{ANCOVA.} Analysis of covariance on $\Phi_\text{after}$ with
$\Phi_\text{before}$ as a covariate and mode as a categorical factor.
This addresses the signal-to-noise problem of $\Delta\Phi$ (the difference of
two noisy measurements) by modeling the baseline statistically rather than
arithmetically.

We apply Bonferroni correction across the nine pairwise comparisons
(adjusted $\alpha = 0.05/9 = 0.0056$). Sensitivity analysis of
$\Phi_\text{approx}$ to weight choices uses Dirichlet concentration near
default weights; robustness is assessed via coefficient of variation
(standard deviation / mean).

%══════════════════════════════════════════════════════════════════════════════
% 3. Results
%══════════════════════════════════════════════════════════════════════════════

\section{Results}

All nine modes completed 100 observation cycles each (900 total). No
exceptions or data exclusions occurred. During development, two
implementation issues were identified and corrected before the results
presented here were obtained (see Section~4.4 for details): (1) coalition
formation used text-token Jaccard similarity that was ineffective for neural
activation summaries, and (2) effective information accumulated linearly with
cycle count due to expanding TPM history. Both issues are fully resolved in
the current data.
Full per-cycle data are available in the supplementary JSONL files.

\begin{table}[h]
\caption{Per-mode $\Phi$ statistics (mean $\pm$ SD over 100 cycles). $\Phi_\text{within}$
and $\Phi_\text{between}$ are from the decomposition (Eq.~\ref{eq:decompose}).
Note that $\Phi_\text{within}$ and $\Phi_\text{after}$ are computed via
different algorithmic paths and have different scales. Complexity is the
activation differentiation metric (standard deviation of processor outputs).}
\label{tab:results}
\centering
\small
\begin{tabular}{lcccccc}
\toprule
Mode & $\Phi_\text{after}$ & $\Phi_\text{before}$ & $\Delta\Phi$ &
$\Phi_\text{within}$ & $\Phi_\text{between}$ & Complexity \\
\midrule
Collaborative           & 1.049 $\pm$ 0.194 & 1.039 & 0.010 & 2.822 & 0.055 & 0.207 \\
Competitive             & 1.046 $\pm$ 0.189 & 1.035 & 0.011 & 2.784 & 0.052 & 0.208 \\
Random                  & 1.049 $\pm$ 0.193 & 1.039 & 0.010 & 2.795 & 0.055 & 0.221 \\
Hybrid                  & 1.047 $\pm$ 0.192 & 1.036 & 0.010 & 2.783 & 0.052 & 0.209 \\
Adaptive                & 1.049 $\pm$ 0.194 & 1.039 & 0.011 & 2.809 & 0.055 & 0.211 \\
No-broadcast            & 1.046 $\pm$ 0.192 & 1.036 & 0.010 & 2.747 & 0.052 & 0.206 \\
Homogeneous competitive & 1.022 $\pm$ 0.189 & 0.984 & 0.039 & 2.790 & 0.028 & 0.113 \\
Single-processor        & 0.471 $\pm$ 0.043 & 0.461 & 0.010 & 1.197 & 0.051 & 0.204 \\
Single self-broadcast   & 0.471 $\pm$ 0.043 & 0.461 & 0.010 & 1.197 & 0.051 & 0.204 \\
\bottomrule
\end{tabular}
\end{table}

\subsection{Finding 1: $\Phi_\text{approx}$ Is Dominated by Processor Count,
Not by Broadcast Architecture}

Table~\ref{tab:results} reveals a pattern fundamentally different from our
pre-registered expectations. All seven multi-processor modes (competitive,
random, no-broadcast, collaborative, hybrid, adaptive, and homogeneous
competitive) produce nearly identical $\Phi_\text{after}$ values in the
range 1.022--1.047. The broadcast mechanism---whether competitive, random,
coalition-based, or absent entirely---has negligible effect on the resulting
$\Phi$. A Mann--Whitney $U$ test comparing competitive versus no-broadcast
(H5) on $\Phi_\text{after}$ yields $d = 0.003$, $p = 0.48$---no significant
difference.

The ANCOVA on $\Phi_\text{after}$ with $\Phi_\text{before}$ as covariate
yields $F(8, 890) = 3257$, $p \approx 0$, but the effect is driven entirely
by two conditions: the single-processor modes ($\Phi_\text{adjusted} \approx
0.95$) and the no-broadcast mode's anomalous baseline. All multi-processor
modes have adjusted means within 0.02 of each other (1.064--1.083). Mode
does not meaningfully differentiate $\Phi$ among multi-processor architectures.

A clustered analysis aggregating the 5 serial cycles per stimulus into
20 independent stimulus-level means per mode corroborates this: competitive
versus no-broadcast yields $d = 0.003$, $p = 0.44$ (Mann--Whitney $U$ on
per-stimulus means).

The $\Phi$ decomposition explains why. $\Phi_\text{within}$ (the effective
information from processor-internal TPMs) is 2.75--2.81 for all multi-processor
modes---including no-broadcast, which now records processor states in
workspace history for comparable TPM construction. $\Phi_\text{between}$
(the MI integration component) is a near-constant 0.051--0.053 across all
modes. $\Phi_\text{approx}$ is dominated by the number of recurrent neurons
contributing to the TPM (5 $\times$ 512 = 2560 for multi-processor modes,
1 $\times$ 512 = 512 for single-processor), not by the broadcast mechanism.

\subsection{Finding 2: Processor Differentiation Modestly but Robustly
Increases $\Phi$}

The one consistent difference among multi-processor modes is between
differentiated and homogeneous processors. Competitive selection with
five structurally differentiated processors yields $\Phi_\text{after} =
1.046 \pm 0.189$, while the same mechanism with five identical Integrator
RNNs yields $1.022 \pm 0.189$ (H7: clustered Mann--Whitney $U$ on
per-stimulus means, $d = 0.13$, $p < 0.001$).

This is the only hypothesis that survives correction across the pairwise
comparisons: processor differentiation provides a small but statistically
reliable $\Phi$ benefit. The homogeneous condition's lower $\Phi$ is
accompanied by a sharply reduced complexity metric (0.113 vs.\ 0.206--0.208),
providing convergent evidence that structural diversity among processors
matters for the system's information-theoretic properties. However, the
effect size is modest ($d = 0.13$), suggesting that differentiation is a
real but minor contributor relative to the dominant effect of total neuron
count.

\subsection{Finding 3: Coalition Formation Is Achieved but Confers No
$\Phi$ Advantage}

After correcting the pairwise agreement mechanism to use neural cosine
similarity rather than text-token Jaccard, coalition modes now form genuine
multi-member coalitions. In collaborative mode, the attention weight standard
deviation increased from 0.0015 to 0.016, with an effective ensemble size
of approximately 2.5 processors (inverse participation ratio of 1/0.4),
indicating that 2--3 of the 5 processors now contribute meaningfully to the
merged representation. The merge differentiation retention score rose from
0.009 to 0.035, confirming that the merge now preserves more of the
individual processors' distinct contributions.

However, this genuine coalition formation does not translate into higher
$\Phi$. Collaborative ($\Phi = 1.047$) and competitive ($\Phi = 1.046$)
are statistically indistinguishable (H1: $d = -0.01$, $p = 0.42$). Coalition
consensus broadcast, even when operating correctly, does not produce higher
causal integration than single-winner selection in the current system.

The adaptive mode tells a consistent story. With $\sigma(\alpha) = 0.022$
in the fixed version (up from the artifactually low 0.0007), the
differentiation gate triggers coalition in a subset of cycles where
processor disagreement exceeds the threshold. Yet the resulting
$\Phi_\text{after}$ (1.046) remains identical to all other multi-processor
modes. Adaptive gating correctly modulates \textit{whether} coalition is
used, but coalition does not change $\Phi$ regardless.

\subsection{Finding 4: The Diagnostic Methodology Is the Primary Contribution}

Two implementation issues were discovered during development and corrected
before the results reported here were obtained: (1) the coalition formation
bug, where \texttt{\_pairwise\_agreement} used Jaccard similarity on
short neural text summaries that lacked token overlap, preventing
multi-member coalitions from ever forming; and (2) the EI accumulation
artifact, where effective information was computed on expanding TPM history
windows, causing $\Phi$ to grow linearly with cycle count and creating a
circular broadcast-vs-no-broadcast comparison (no-broadcast had zero history
by construction, yielding $\Phi_\text{within} = 0$).

These issues were identified \textit{because} the diagnostic tools developed
in this work---$\Phi$ decomposition, per-cycle tracking, and attention
weight monitoring---made them visible. The $\Phi_\text{within}$ /
$\Phi_\text{between}$ decomposition revealed the accumulation artifact (a
flat $\Phi_\text{between}$ coupled with a linearly growing
$\Phi_\text{within}$). Per-cycle coalition size tracking revealed that
92\% of ``collaborative'' cycles had a coalition of size 1. The
methodological contribution of this work is therefore not the specific
empirical findings---which were substantially revised by these corrections---but
the diagnostic framework that enabled their discovery.

\subsection{Processor Functional Similarity}

A critical auxiliary finding concerns the relationship between structural
and functional processor differentiation. Pairwise cosine similarity between
processor activation vectors (pre-stimulus resting state, 512-dimensional)
averages 0.027 across the 10 processor pairs, confirming that the five
$W_\text{rec}$ initialization strategies produce structurally distinct
dynamical regimes. However, the similarity distribution is skewed: 9 of 10
pairs fall below 0.07 (range: $-$0.069 to 0.052), while the
reasoner--integrator pair reaches 0.20.

This structural differentiation does not translate into functional
differentiation when processors process the same stimulus. In collaborative
mode, the attention-weighted merge produces a mean cosine similarity of 0.92
between processor outputs and the workspace vector (standard deviation 0.02),
indicating that despite structurally distinct recurrent weights, the five
processors converge to highly similar activation patterns when driven by a
shared stimulus. The effective ensemble size of the coalition merge is
approximately 2.5 processors (inverse participation ratio), consistent with
2--3 processors contributing meaningfully while the remainder produce
redundant representations.

This dissociation between structural and functional differentiation explains
why coalition formation, even when operating correctly, confers no $\Phi$
advantage: the processors agree too much for coalition to add information
beyond what any single processor already provides. The implication for
multi-agent architectures is that \textit{structural} heterogeneity
(different model architectures or connectivity patterns) is insufficient to
guarantee \textit{functional} heterogeneity---the stimulus itself drives
convergence.

\subsection{Calibration}

A calibration on three reference systems matched to the experiment scale
($n_\text{neurons}=512$, $n_\text{processors}=5$, $n_\text{clusters}=3$,
$n_\text{history}=10$) provides interpretive bounds: random noise yields
$\Phi_\text{approx} = 0.738$, deterministic sine dynamics 0.712, and
correlated signal 0.580. These values are below the experimental
multi-processor $\Phi$ range (1.02--1.05), consistent with the structured
recurrent dynamics ($W_\text{rec}$ with spectral radius $< 0.9$) producing
stronger causal constraints than unstructured or deterministic reference
systems. The calibration confirms that $\Phi_\text{approx}$ is sensitive to
genuine causal structure (RNN $>$ noise), but the narrow spread of
calibration values (0.58--0.74 vs.\ the experimental range of 1.02--1.05)
indicates that recurrent connectivity is the dominant contributor to the
metric.

%══════════════════════════════════════════════════════════════════════════════
% 4. Discussion
%══════════════════════════════════════════════════════════════════════════════

\section{Discussion}

\subsection{The Dominance of Neuron Count over Broadcast Architecture}

The most important empirical result is a null result: among multi-processor
modes, broadcast mechanism explains essentially none of the variance in
$\Phi_\text{approx}$. Whether processors compete (competitive), are chosen
at random (random), form coalitions (collaborative, hybrid), adaptively
gate themselves (adaptive), or do not broadcast at all (no-broadcast),
the resulting $\Phi$ is the same within a band of approximately 0.02.
The ANCOVA's large $F$-statistic is driven by the single-processor condition
alone; among the seven multi-processor modes, adjusted means are nearly
identical.

This finding has a straightforward interpretation: $\Phi_\text{approx}$
as implemented here is dominated by the number of recurrent neurons
contributing to the state transition probability matrix. Five processors
with 512 neurons each produce approximately twice the $\Phi$ of one such
processor (1.05 vs.\ 0.47), consistent with the factor of five in neuron
count. The recurrent weights $W_\text{rec}$ of individual processors
generate the causal structure that EI measures; adding more processors
adds more causal structure, nearly additively, regardless of how they
communicate.

This is \textit{not} a failure of the experimental design. It is a finding
about the nature of $\Phi$ in feedforward-recurrent architectures: the
effective information arises primarily from within-processor recurrence,
and inter-processor communication contributes negligibly at the level of
the current $\Phi$ approximation. The $\Phi_\text{between}$ component
(MI integration) is essentially flat across all modes (0.051--0.053),
confirming that the current MI estimation pipeline does not capture
broadcast-mediated information flow.

\subsection{Processor Differentiation: A Small but Reliable Effect}

The one signal that survives all corrections and analysis methods is the
differentiation effect. Structurally heterogeneous processors
(competitive: $\Phi = 1.046$) modestly but reliably outperform structurally
homogeneous processors (homogeneous competitive: $\Phi = 1.022$). At the
stimulus level (20 independent means per mode), this difference is
statistically significant ($d = 0.13$, $p < 0.001$).

The effect size is small because differentiation only affects the
complexity term ($\omega_\text{diff} = 0.25$) of the composite $\Phi$,
and because the homogeneous processors still differ in their random seeds
and noise realizations even though they share the same $W_\text{rec}$
topology. A maximal differentiation manipulation---five randomly initialized
unstructured $W_\text{rec}$ matrices---would likely produce a larger effect
but at the cost of interpretability (no systematic specialization).

For multi-agent AI, this finding carries a tempered message: agent
heterogeneity improves information integration, but the improvement is
marginal unless coupled with mechanisms that amplify differentiation in
the representations that agents actually produce when processing shared
inputs.

\subsection{Coalition Formation: Achieved but Inconsequential}

After correcting the pairwise agreement mechanism, coalition modes now
form genuine multi-member coalitions. The effective ensemble size increased
from approximately 1.0 to approximately 2.5 processors, and merge
differentiation retention improved by a factor of 4. The attention weight
standard deviation rose from an artifactually low 0.0015 to a more realistic
0.016, indicating that the neural cosine similarity measure is capturing
meaningful variation among processor outputs.

Yet this operational success does not translate into any $\Phi$ advantage.
Coalition modes produce the same $\Phi$ as competitive selection because
the merged representation---whether formed by mean, attention-weighted
average, or concatenation+projection---enters the workspace as a single
activation vector. The TPM constructed from workspace history does not
distinguish between ``this vector came from one processor'' and ``this
vector came from a coalition of three.'' The causal structure downstream
of the workspace is identical.

This finding suggests that coalition consensus broadcast and competitive
winner-take-all broadcast are \textit{causally equivalent} at the level of
the global workspace state---a result that constrains the scope of CGWT's
claims. If coalition is to produce higher $\Phi$, the benefit must arise
either upstream of the workspace (in the selection process itself) or
through a representation that preserves coalition structure within the
broadcast (e.g., concatenation without projection, or a multi-modal
workspace).

\subsection{Methodological Contribution: Diagnostic Tools That Found Real Bugs}

The two implementation issues corrected during development were discovered
\textit{because} the diagnostic tools developed as part of this work made
them visible. The $\Phi_\text{within}$ / $\Phi_\text{between}$ decomposition
revealed an anomalous pattern---flat $\Phi_\text{between}$ coupled with
linearly growing $\Phi_\text{within}$---that pointed to the EI accumulation
artifact. Per-cycle coalition size tracking revealed that 92\% of
``collaborative'' cycles had a single-member coalition, exposing the
pairwise agreement bug. Per-cycle $\Phi$ tracking enabled the accumulation
artifact to be quantified precisely (linear slope of $\sim$0.33 per cycle).

These diagnostic capabilities---decomposing $\Phi$ into within- and
between-processor components, tracking attention weight distributions,
monitoring coalition sizes per cycle, and computing merge fidelity
metrics---constitute a transferable toolkit. Any computational study
of broadcast-mediated information integration would benefit from deploying
these diagnostics, regardless of the specific theoretical framework.

\subsection{Relation to the MaxCal IIT--FEP Bridge}

The null result for broadcast mechanism effects does not diminish the
relevance of the MaxCal bridge~\cite{kearney2026information}. The coalition
selection objective (Eq.~\ref{eq:maxcal_coalition}) remains valid as a
formalization, but the current implementation only operationalizes the
first term (cosine similarity to workspace prior). The second term---the
prediction error penalty---is absent, and the finding that all broadcast
mechanisms are causally equivalent at the workspace level suggests that
the missing term is where the differentiation would manifest. Future
implementations should operationalize the full objective, using
processor-specific prediction errors to weight coalition contributions
such that processors with lower prediction error (better world model
alignment) receive higher weight, and coalitions that minimize collective
prediction error are preferentially selected.

\subsection{Limitations}

\begin{enumerate}
\item \textbf{$\Phi$ approximation validity.} Our $\Phi_\text{approx}$ has
not been validated against exact $\Phi$ on small tractable systems. This
limitation is inherent to IIT itself---exact $\Phi$ is computable only for
systems with $N \leq 16$ binary units. Our $\Phi_\text{approx}$ uses four
documented simplifications (Section~2.3), each with a known bias direction.
If the approximation is valid, we predict that in a tractable binary system
($N = 10$, $M = 2$--$3$ processors), $\Phi_\text{approx}$ and true $\Phi$
should exhibit a Spearman rank correlation $\geq 0.8$ across conditions with
varying causal structure. This validation remains the highest-priority future
work; until completed, all findings based on $\Phi_\text{approx}$ rank-ordering
should be interpreted as contingent on this predicted correlation.

\item \textbf{TPM window size and cycle count.} The fixed TPM window
( = 10$) and the small number of cycles per stimulus (5) mean that
$\Phi$ is measured during transient dynamics rather than at steady state.
Extending to 10--15 cycles per stimulus with a proportionally larger TPM
window would distinguish transient from stable $\Phi$ values and may
reveal broadcast-dependent effects that emerge only after the system
settles.

\item \textbf{Coalition merge causal equivalence.} The finding that all
broadcast mechanisms produce equivalent $\Phi$ is partly a consequence of
the TPM operating on workspace states: once a vector enters the workspace,
its provenance (single processor vs.\ coalition merge) is lost. A causal
analysis that preserves coalition membership structure through the
broadcast may recover broadcast-dependent $\Phi$ differences.

\item \textbf{Processor differentiation.} Structural $W_\text{rec}$
differences produce functionally similar activations under shared stimulus
input ($\cos\theta > 0.92$ in collaborative and adaptive modes). More
aggressive specialization---different input projections, heterogeneous
timescales ($T \in \{15, 20, 25, 30, 35\}$), or non-tanh activation
functions---may be required for coalition to confer a $\Phi$ benefit.

\item \textbf{Stimulus encoding.} Natural language stimuli are encoded into
32-dimensional vectors via deterministic hashing, discarding most semantic
content. A stimulus ablation (replacing stimulus vectors with random noise)
is needed to measure the stimulus-dependence of $\Phi$.

\item \textbf{Statistical dependence.} The 5 serial cycles per stimulus
are not independent (each cycle's output becomes the next cycle's context).
We report stimulus-level clustered analyses (20 independent means per mode)
alongside per-cycle tests; the clustered results are more conservative and
should be preferred for hypothesis testing.
\end{enumerate}

\subsection{Downstream Implications}

Despite the null results for broadcast mechanism effects, several findings
carry implications for adjacent fields:

\textbf{Multi-agent AI.} The differentiation-prerequisite finding (H7:
 = 0.13$,  < 0.001$) corroborates the intuition that agent heterogeneity
improves system-level properties, but the small effect size tempers
practical significance. Agent diversity alone, without mechanisms that
amplify representational differences, yields marginal gains. The coalition
formation mechanism (neural cosine similarity for member selection)
demonstrates a concrete alternative to text-based consensus in LLM
multi-agent debate frameworks.

\textbf{Neuromorphic computing.} The approximate linear scaling of $\Phi$
with total neuron count (0.47 for 512 neurons, 1.05 for 2560 neurons)
provides an empirical benchmark for $\Phi_\text{approx}$ as a function of
system size on recurrent architectures. The differentiation-prerequisite
finding translates to a testable prediction for heterogeneous vs.homogeneous core configurations on neuromorphic hardware.

\textbf{Computational consciousness methodology.} The diagnostic
toolkit---$\Phi$ decomposition, attention weight tracking, coalition size
monitoring, per-cycle $\Phi$ trajectories---proved its value by identifying
two non-trivial implementation issues during development (the EI
accumulation artifact and the coalition formation bug). These tools
constitute a transferable contribution to the methodology of computational
consciousness science, applicable to any study computing $\Phi$ on
multi-component systems.
\end{enumerate}

\subsection{Future Work}

\textbf{1. Exact $\Phi$ validation.} Validating $\Phi_\text{approx}$ against
exact $\Phi$ on small tractable systems ($N = 6$--$12$ binary units) is
the highest-priority next step. If the rank correlation is $\geq 0.8$, the
current findings gain methodological credibility.

\textbf{2. Coalition-aware causal analysis.} Modifying the TPM construction
to preserve coalition membership structure through the broadcast---so that
a coalition-merged workspace state is distinguishable from a single-winner
state---may recover broadcast-dependent $\Phi$ differences.

\textbf{3. MaxCal coalition optimization.} Implementing the full coalition
selection objective (Eq.~\ref{eq:maxcal_coalition}) with the prediction
error penalty term would test whether optimization-based coalition selection
can produce a $\Phi$ advantage over simple scoring.

\textbf{4. Stimulus ablation and extended cycles.} Replacing stimulus
vectors with random noise would measure stimulus-dependence. Extending
cycles per stimulus to 10--15 would allow $\Phi$ to reach steady state.

\textbf{5. Cross-architecture replication.} Testing whether $\Phi$ continues
to scale approximately linearly with neuron count on LSTM and GRU
architectures would establish generality.

\textbf{6. Processor specialization beyond connectivity.} More aggressive
specialization---different input projections, heterogeneous timescales,
non-tanh activations---should be explored to test whether sufficient
functional differentiation enables coalition to confer a $\Phi$ benefit.

\section{Conclusions}

We implemented a 5-processor recurrent neural network system with global
workspace broadcast and computed a tractable $\Phi$ approximation from
neural state transition probability matrices across nine experimental
modes (900 observation cycles). Two implementation issues---an EI
accumulation artifact from expanding TPM windows and a coalition formation
bug from mismatched agreement metrics---were identified through the
diagnostic toolkit developed in this work and corrected during development.

After corrections, the principal empirical finding is a null result:
\textbf{among multi-processor modes, broadcast mechanism explains
essentially none of the variance in $\Phi_\text{approx}$}. All seven
multi-processor modes produce $\Phi$ in the narrow range 1.022--1.047,
dominated by the total recurrent neuron count (2560 for five processors
vs.\ 512 for one). The one reliable signal is a small differentiation
effect: structurally heterogeneous processors modestly outperform
homogeneous ones (clustered $d = 0.13$, $p < 0.001$).

The CGWT hypothesis that coalition consensus broadcast produces higher
$\Phi$ than competitive selection is not supported. After correcting the
coalition formation mechanism, genuine multi-member coalitions now form
(effective ensemble size $\approx 2.5$), but the merged representation
produces causally equivalent workspace states to single-winner selection.

The contribution of this work is therefore methodological rather than
confirmatory. The diagnostic framework---$\Phi$ decomposition into within-
and between-processor components, attention weight distribution tracking,
coalition size monitoring, per-cycle $\Phi$ trajectory analysis, and merge
fidelity metrics---enabled the discovery of two non-trivial implementation
issues and provides a transferable toolkit for computational studies of
broadcast-mediated information integration. The neuron-count dominance
finding constrains the scope of claims that broadcast architecture alone
can produce measurable $\Phi$ differences under the current approximation,
and the differentiation-prerequisite finding provides a modest but reliable
empirical anchor for processor specialization as a design principle.

\vspace{6pt}
\noindent\textbf{Data and Code Availability:} All code, experiment
configurations, raw data (per-cycle JSONL), and analysis scripts are
available at \url{https://github.com/guyu-adam/noesis} under the MIT license.
The specific experiment data reported in this paper correspond to commit
\texttt{ef5a1ff} and the configuration file
\texttt{experiment\_config\_20260521\_142426.json}.

\vspace{6pt}
\noindent\textbf{Acknowledgments:} The author thanks the developers of
PyTorch, CuPy, NumPy, and SciPy for the open-source tools that made this
work possible. The development of the diagnostic toolkit benefited
substantially from systematic internal validation and debugging, during
which two non-trivial implementation issues were identified and corrected.

%══════════════════════════════════════════════════════════════════════════════
% References
%══════════════════════════════════════════════════════════════════════════════

\begin{thebibliography}{99}

\bibitem{baars1988cognitive}
Baars, B.J. \textit{A Cognitive Theory of Consciousness}; Cambridge University Press: Cambridge, UK, 1988.

\bibitem{dehaene2011experimental}
Dehaene, S.; Changeux, J.P. Experimental and theoretical approaches to conscious processing. \textit{Neuron} \textbf{2011}, \textit{70}, 200--227.

\bibitem{tononi2004information}
Tononi, G. An information integration theory of consciousness. \textit{BMC Neurosci.} \textbf{2004}, \textit{5}, 42.

\bibitem{tononi2016integrated}
Tononi, G.; Boly, M.; Massimini, M.; Koch, C. Integrated information theory: from consciousness to its physical substrate. \textit{Nat. Rev. Neurosci.} \textbf{2016}, \textit{17}, 450--461.

\bibitem{ferrante2025adversarial}
Ferrante, O.; Gorska-Klimowska, U.; Henin, S.; Hirschhorn, R.; Khalaf, A.; Lepauvre, A.; Liu, L.; Richter, D.; Vidal, Y.; Bonacchi, N.; et al. Adversarial testing of global neuronal workspace and integrated information theories of consciousness. \textit{Nature} \textbf{2025}, \textit{642}, 133--142.

\bibitem{luppi2024synergistic}
Luppi, A.I.; Mediano, P.A.M.; Rosas, F.E.; Allanson, J.; Pickard, J.D.; Carhart-Harris, R.L.; Williams, G.B.; Craig, M.M.; Finoia, P.; Owen, A.M.; et al. A synergistic workspace for human consciousness revealed by integrated information decomposition. \textit{eLife} \textbf{2024}, \textit{12}, e88173.

\bibitem{kearney2026information}
Kearney, A. Information as Maximum-Caliber Deviation. \textit{arXiv} \textbf{2026}, 2605.12536.

\bibitem{barrett2026iit}
Barrett, A.B.; Milinkovic, B.; Mediano, P.A.M.; Rosas, F.E.; Bor, D.; Barnett, L.; Seth, A.K. IIT: The Good, the Bad and the Misunderstood. \textit{arXiv} \textbf{2026}, 2604.11482.

\bibitem{iwmt2020}
Safron, A. An Integrated World Modeling Theory (IWMT) of consciousness: combining integrated information and global neuronal workspace theories with the free energy principle and active inference framework; toward solving the hard problem and characterizing agentic causation. \textit{Front. Artif. Intell.} \textbf{2020}, \textit{3}, 30.

\bibitem{safron2020integrated}
Safron, A. Integrated world modeling theory expanded: implications for the future of consciousness. \textit{Front. Comput. Neurosci.} \textbf{2022}, \textit{16}, 1--18.

\bibitem{godelos}
Hirst, O.C. (Steake). GodelOS: A Consciousness Operating System for LLMs. GitHub repository, 2025--2026. \url{https://github.com/steake/GodelOS}.

\bibitem{multiagent_coordination}
Du, Y.; Li, S.; Torralba, A.; Tenenbaum, J.; Mordatch, I. Improving Factuality and Reasoning in Language Models through Multiagent Debate. \textit{arXiv} \textbf{2024}, 2305.14325. [Representative example.]

\bibitem{hoel2013causal}
Hoel, E.P.; Albantakis, L.; Tononi, G. Quantifying Causal Emergence Shows That Macro Can Beat Micro. \textit{Proc. Natl. Acad. Sci. USA} \textbf{2013}, \textit{110}, 19790--19795.

\bibitem{phua2025}
Phua, Y.J. Can We Test Consciousness Theories on AI? Ablations, Markers, and Robustness. \textit{arXiv} \textbf{2025}, 2512.19155.

\bibitem{tononi2015}
Tononi, G.; Koch, C. Consciousness: here, there and everywhere? \textit{Philos. Trans. R. Soc. B} \textbf{2015}, \textit{370}, 20140167.

\bibitem{friston2010}
Friston, K. The free-energy principle: a unified brain theory? \textit{Nat. Rev. Neurosci.} \textbf{2010}, \textit{11}, 127--138.

\bibitem{kriegeskorte2008representational}
Kriegeskorte, N.; Mur, M.; Bandettini, P. Representational similarity analysis---connecting the branches of systems neuroscience. \textit{Front. Syst. Neurosci.} \textbf{2008}, \textit{2}, 4.

\bibitem{mediano2021taxonomy}
Mediano, P.A.M.; Rosas, F.E.; Luppi, A.I.; Carhart-Harris, R.L.; Bor, D.; Seth, A.K.; Barrett, A.B. Towards an Extended Taxonomy of Information Dynamics via Integrated Information Decomposition. \textit{arXiv} \textbf{2021}, 2109.13186.

\bibitem{oizumi2014}
Oizumi, M.; Albantakis, L.; Tononi, G. From the phenomenology to the mechanisms of consciousness: integrated information theory 3.0. \textit{PLoS Comput. Biol.} \textbf{2014}, \textit{10}, e1003588.

\end{thebibliography}

\end{article}
